
\documentclass[12pt,a4paper]{report}

\usepackage[a4paper,margin=1in]{geometry}
\usepackage{graphicx}
\usepackage{float}
\usepackage{amsmath}
\usepackage{booktabs}
\usepackage{longtable}
\usepackage{array}
\usepackage{hyperref}
\usepackage{listings}
\usepackage{xcolor}
\usepackage{fancyhdr}
\usepackage{titlesec}

% Rename bibliography heading
\renewcommand{\bibname}{References}

\pagestyle{fancy}
\fancyhf{}
\rhead{CS3243 Project Report}
\lhead{IoT Devices and Applications}
\cfoot{\thepage}
\setlength{\headheight}{14.5pt}

\fancypagestyle{plain}{
    \fancyhf{}
    \rhead{CS3243 Project Report}
    \lhead{IoT Devices and Applications}
    \cfoot{\thepage}
}

\titleformat{\chapter}[display]
    {\normalfont\bfseries\Huge}
    {\chaptername\ \thechapter}
    {0pt}
    {\Huge}
\titlespacing*{\chapter}{0pt}{0pt}{24pt}
\titleformat{\section}{\large\bfseries}{\thesection}{1em}{}
\titleformat{\subsection}{\normalsize\bfseries}{\thesubsection}{1em}{}

\title{
    \textbf{Hospital Automated Chemical Dosing System}\\
    \textbf{Design Report} \\
    \large CS3243 - IoT Devices and Applications
}

\author{
    KR Kanes Rakeshan \\
    230518C
}

\date{\today}

\begin{document}

\maketitle
\tableofcontents

\newpage

\chapter*{Introduction}
\addcontentsline{toc}{chapter}{Introduction}

\section*{Background}
Hospitals use specialized disinfectant chemicals to sterilize medical instruments and equipment.

\section*{Objectives}
\begin{itemize}
    \item Design a complete IoT-based chemical dispensing solution
    \item Implement automatic refill control using ESP32
    \item Monitor water level using non-contact sensing
    \item Implement fault detection and alarm generation
\end{itemize}

\section*{Proposed IoT System Architecture}
\subsection*{System Overview}
The proposed system consists of:
\begin{itemize}
    \item Purified water tank
    \item Concentrated chemical supply tank
    \item Mixing chamber
    \item Dispensing tanks
    \item ESP32-based controller
\end{itemize}

\chapter{Method to Measure Water Level and Remaining Chemical Levels}

\section{Water Level Measurement and Chemical Estimation}

\subsection{Water Tank Measurement}
An ultrasonic sensor was selected for non-contact level measurement.

\begin{equation}
Water\ Level = Tank\ Height - Measured\ Distance
\end{equation}

\subsection{Chemical Estimation}
\begin{equation}
Remaining\ Chemical = Initial\ Volume - \sum Dispensed\ Volume
\end{equation}

\section{Notes / Remaining Work}
\noindent\textit{Placeholder. Add constraints and justify sensor/method choices based on the given constraints.}

\chapter{Component Selection (Valves, Flow Meters, Pump, Sensors, etc.)}

\section{Component Selection}

\subsection{ESP32 Controller}
ESP32 was selected due to:
\begin{itemize}
    \item Built-in WiFi
    \item PWM support
    \item ADC support
    \item Low cost
\end{itemize}

\section{Market Evaluation (Minimum Three Products)}
\noindent\textit{Placeholder. Compare at least three products for each required component (features + cost).}

\chapter{Failure Identification, Detection, and Fail-Safe Methods}

\section{Fault Detection}
\begin{table}[H]
\centering
\begin{tabular}{lll}
\toprule
Failure & Detection & Action \\
\midrule
Pump Failure & No flow pulses & Alarm \\
Blocked Valve & No flow & Stop refill \\
Leakage & Flow mismatch & Shutdown \\
\bottomrule
\end{tabular}
\caption{Fault Detection Mechanisms}
\end{table}

\section{Sensor-Based and Computational Failure Detection}
\noindent\textit{Placeholder. List potential failures, detection method(s), and fail-safe actions.}

\chapter{Data Collection and Communication to Central Controller}
\noindent\textit{Placeholder. Evaluate at least two communication methods suitable for the distance/distribution.}

\chapter{Complete System Schematic Block Diagram}
\noindent\textit{Placeholder. Insert the complete schematic block diagram of all system components.}

\chapter{Controller and Sensor Connection Schematic}
\noindent\textit{Placeholder. Insert a schematic of the controller and connected sensors.}

\chapter{Operational State Diagram / Flow Chart for IoT Components}
\noindent\textit{Placeholder. Insert the algorithm flow chart / state diagram per IoT component.}

\chapter{Simulation Implementation (TinkerCad/WokWi) and Source Code Attachment}

\section{Simulation}
The system was simulated using WokWi with:
\begin{itemize}
    \item ESP32
    \item Ultrasonic sensor
    \item Flow pulse simulation
    \item Alarm system
\end{itemize}

\section{Source Code (Attachment)}
\noindent\textit{Placeholder. Attach well documented source code and reference it here.}

\chapter*{Conclusion}
\addcontentsline{toc}{chapter}{Conclusion}
An IoT-based automated hospital chemical dispensing system was successfully designed and simulated.

\begin{thebibliography}{9}
\addcontentsline{toc}{chapter}{References}

% \bibitem{book1}
% Adrian McEwen and Hakim Cassimally,
% \textit{Designing the Internet of Things}.

% \bibitem{esp32}
% Espressif Systems,
% ESP32 Technical Reference Manual.

\end{thebibliography}

\end{document}
